\begin{abstract}
Blood vessels are the main reason why nodule detectors in chest CT look in the wrong place: where a vessel bends or branches it produces a bright, compact spot that resembles a small lesion. Predicting the vessels from the surrounding tissue and removing them before the search is an appealing remedy, but it has been unclear whether such a vessel forecast helps because it is in the right place or merely because it gives the detector more information. We answer this question with a controlled experiment in which each lesion-free image region and its copy with an inserted lesion share the same surroundings, and therefore the same vessel forecast, while a rotated copy of the forecast serves as a misplaced control. The answer is that placement matters, but mainly for where the detector looks rather than for how many regions it accepts. At the same false-alarm rate, a correctly placed forecast moves the detector's mark onto the lesion for about 6 additional insertions per 100 on held-out LUNA16 scans, and by about 3 per 100 on scans with expert-annotated native vessels, whereas the number of accepted regions barely changes. Simply subtracting the forecast works as well as more flexible models of the vessel. These results suggest that vessel-aware detectors should be evaluated by where they point, not only by what they accept.
\end{abstract}
